\newpage
\section{Phénomènes aléatoires}
Les probabilités sont un outil que l'on utilise pour étudier des phénomènes particuliers : ceux dont on ne peut prédire avec certitude les résultats.

\ex{\begin{eleve}
    Un lancer de dé.
\end{eleve}}


\defn{Une \textbf{épreuve} du phénomène est un résultat possible du phénomène.\\
Notation : $\omega$.}\\

\ex{
Une épreuve du lancer de dé est \begin{eleve} la valeur 5.
\end{eleve}}
\defn{La \textbf{catégorie d'épreuves} est l'ensemble de toutes les épreuves du phénomène.\\
Notation : $\Omega$.}\\

\ex{
La catégorie d'épreuves du lancer de dé est \begin{eleve} l'ensemble $\{1;2;3;4;5;6\}$\end{eleve}}

\defn{Un \textbf{évènement} est un ensemble d'épreuves. C'est donc un sous-ensemble de la catégorie d'épreuves.\\
Notation : lettre majuscule $A, B ...$}.\\

\ex{
On peut avoir plusieurs évènements différents au lancer de dé, dépendant de ce qui nous intéresse.\\
Si on s'intéresse aux lancers de dés à résultats impairs, l'évènement $A$ est \begin{eleve}$A = \{1;3;5\}$.\end{eleve}\\
Si on s'intéresse aux lancers de dés à résultats strictement inférieurs à 5, l'évènement $B$ est \begin{eleve}$B = \{1;2;3;4\}$.\end{eleve}\\
Si on s'intéresse aux lancers de dés à résultat égal à 2, l'évènement $C$ est \begin{eleve}$C = \{2\}$.\end{eleve}}\\

Ces évènements sont des ensembles, on peut donc utiliser les opérations que l'on a définies précédemment sur ces évènements.
\ex{.\\
L'évènement $A$ est l'ensemble des résultats impairs au lancer de dés.\\
L'évènement $B$ est l'ensemble des résultats strictement inférieurs à 5.\\
L'évènement $C$ est l'ensemble des résultats égaux à 2.\\
L'évènement $A^c$ est \begin{eleve}l'ensemble des résultats pairs au lancer de dés.\end{eleve}\\
$A^c = $\begin{eleve}$\Omega \setminus A = \{1;2;3;4;5;6\} \setminus \{1;3;5\} = \{2;4;6\}$.\end{eleve}\\
L'évènement $B^c$ est \begin{eleve}
    l'ensemble des résultats supérieurs ou égaux à 5.\\
\end{eleve}
$B^c = $\begin{eleve}$\Omega \setminus B = \{1;2;3;4;5;6\} \setminus \{1;2,3,4\} = \{5;6\}$.\end{eleve}\\
L'évènement $C^c$ est \begin{eleve}l'ensemble des résultats différents de 2\end{eleve}.\\
$C^c = $\begin{eleve}$\Omega \setminus C = \{1;2;3;4;5;6\} \setminus \{2\} = \{1;3;4;5;6\}$.\end{eleve}\\
$A\cap B = $\begin{eleve}
 $\{1;3\}$
\end{eleve}
\hfill
$A\cap C = $\begin{eleve}
 $\emptyset$ (A et C sont disjoints)
\end{eleve}
\hfill
$A\cup B = $\begin{eleve}
 $\{1;2;3;4;5\}$
\end{eleve}